% Keep to one paragraph, ~210 words. Refresh numbers when final endpoints land.
Powerful pretrained video VAEs compress video into rich but \emph{entangled}
latents: object identity, object motion, and camera motion are fused into one code
that cannot be read, edited, or predicted factor by factor. Recovering such
structure normally means training a new VAE from scratch. We present \method{}, a
lightweight adapter that re-encodes the latent of \emph{any frozen} video VAE into
named, independently accessible slots --- a spatial static code (identity), a
per-frame spatial dynamics code (object motion), and a camera code --- with a
decoder structurally prevented from mixing them, and \emph{no} VAE retraining.
Carrying this factorization to real moving-camera robot video exposed a general
failure: a \emph{global} per-frame dynamics vector reconstructs real textured video
$17$\,dB below the frozen-VAE ceiling, because the decoder can only broadcast it
uniformly across space. Giving the dynamics code a spatial axis raises held-out
reconstruction by $+4.4$\,dB, which a matched-rate ablation attributes to spatial
\emph{structure} ($+2.7$\,dB) rather than added rate ($+0.6$\,dB). We show the code
is cleanly \emph{disentangled} (a diagonal-dominant cross-probe), controllable
(identity/motion swap), and camera-aware (a known-pose path for viewpoint stability
that entangled and camera-blind baselines lack), and that its dynamics code
transfers to robot action read-out. We are explicit about what \method{} is
\emph{not}: at extreme compression it does not beat foundation encoders on semantic
accuracy or classical codecs on rate--distortion --- we report both. Its
contribution is a cheap, camera-aware \emph{disentanglement} of frozen video-VAE
latents that entangled predictive VAEs and camera-blind decomposition VAEs do not
provide.
